\section{Erasable is not beneficial (RQ2)}
\label{sec:rq2}

\begin{figure}[t]
\centering
\includegraphics[width=\linewidth]{fig3_erasure_no_benefit.pdf}
\caption{Erasure strength vs target benefit across 40 held-out-target cells. \textbf{A}: no eraser
achieves a proven target-bAcc gain (benefit\_claimable $=0/40$); many cells task-collapse (INLP) or
harm a binary-latent target (LEACE/RLACE on Lee/Cho EEGNet). \textbf{B}: the apparent negative
$\rho(\text{subject-removed},\Delta\text{target bAcc})$ is \emph{not robust} --- it flips to $+0.54$ on
the principled-eraser subset (LEACE/RLACE) and is non-significant when INLP or random-$k$ are dropped;
it is therefore not reported as a finding.}
\label{fig:erasure}
\end{figure}

Subject signal is erasable: LEACE drives linear subject decodability to chance on both backbones (a
nonlinear MLP residual persists), and INLP drives it to chance entirely. The question is whether that
removability confers a target benefit.

\paragraph{No certified benefit.} Deploying source-fitted erasers to held-out target subjects, we find
\textbf{no eraser achieves a proven target-bAcc gain in any of 40 cells} (benefit\_claimable $=0/40$
under a rule requiring a positive point estimate whose 95\% CI lower bound exceeds zero, no
task-collapse, no binary-latent harm, and no random-$k$-matched non-specific NLL move;
Figure~\ref{fig:erasure}A). Erasure can actively harm the target: INLP over-erases and drives task to
chance in several cells, and on the compact binary EEGNet latent (Lee2019, Cho2017) LEACE and RLACE
drive task-bAcc to chance, a task-safety failure. On 2a-TSMNet the small NLL improvement under LEACE is
matched by a same-rank random removal that leaves the subject fully decodable, so that NLL movement is
non-specific regularization, not domain removal.

\paragraph{The negative correlation is not a finding.} Pooled across all 40 cells,
$\rho(\text{subject-removed}, \Delta\text{target bAcc})$ is negative and excludes zero. A sensitivity
battery (Figure~\ref{fig:erasure}B) shows this is not robust: it is non-significant when INLP or
random-$k$ are removed, and it \emph{flips to $+0.54$} (excluding zero) on the principled-eraser
subset (LEACE/RLACE only). The negative sign is driven by INLP over-erasure and the random-$k$ anchor,
not by a ``more removal, worse target'' law. We therefore report only the robust result --- the absence
of a certified benefit --- and do not headline a negative association. The upshot: subject signal is
erasable, but erasure strength does not certify target benefit.
